% =============================================================================
% §1.7 重写草稿（独立片段，未 patch 进主稿）
% 目标：把 §1.7 从"训练目标+候选并置"的 4 小节，重构为真正的"实验设置与比较
%       协议"6 小节；训练目标(loss)留此但重定位为实验设置组件；把递推协议、指标、
%       协议轴与初值扰动（含 Remus100-grounded 噪声预算表）、候选/消融/证据纳入
%       逐项展开，不再一笔带过。
% 红线：三轴正交；噪声 profile 仅为 filtered-state 鲁棒性正则、非导航后验；§1.7 不
%       预锁 §1.8 证据子集；不报告性能数值/排序/"最重要结构"结论。
% 交叉引用沿用方法章既有 label，已逐一核对。
% =============================================================================

\section{实验设置与比较协议}
\label{sec:training-baselines-protocol}

第~\ref{sec:structured-continuous-model} 节给出的结构保持参数化，已在可训练函数类层面保证了正定逆质量、半正定耗散与斜对称零功率耦合等结构约束；第~\ref{sec:energy-power-properties} 节进一步给出了静水机械核心的功率平衡性质。本节规定如何在统一的实验设置下检验这些结构约束在长期递推预测中的作用，包含六个组成部分：可训练参数化与自由度的界定（第~\ref{subsec:structure-preserving-param} 小节）、控制块训练目标（第~\ref{subsec:training-objective} 小节）、长期自由递推验证协议（第~\ref{subsec:rollout-protocol} 小节）、评价指标与内部诊断量（第~\ref{subsec:metrics} 小节）、协议轴与初值扰动设计（第~\ref{subsec:protocol-axes} 小节），以及候选模型、结构消融与证据纳入规则（第~\ref{subsec:candidates-evidence} 小节）。贯穿全节的基本区分是：由参数化直接保证的性质构成模型的\emph{结构约束}，训练只在保结构的函数类内拟合数据驱动的自由度；而由数值积分、状态转换、数据分布与初值扰动共同决定的性质，则通过训练目标、长期递推与诊断指标来\emph{检验}。本节只规定设置与协议，不报告任何性能数值、排序或结构作用结论；实际进入定量比较的证据子集及其结果，留待第~\ref{sec:results-structure-evidence} 节按统一证据口径报告。

\subsection{结构保持参数化与可训练自由度}
\label{subsec:structure-preserving-param}

模型的结构假设通过参数化方式在函数类层面强制，而非通过训练惩罚软性鼓励。具体而言，逆质量映射按式~\eqref{eq:mass-parameterization} 参数化为对称正定矩阵，使相对动量与速度的配对始终良定；非保守广义力按式~\eqref{eq:nonconservative-force} 分解为半正定耗散 \(D_\theta(\xi)\succeq 0\)、斜对称零功率耦合 \(J_\theta(\xi)=-J_\theta(\xi)^\top\) 与可学习广义力端口 \(\tau_\theta\)；保守广义力由标量势能 \(V(q)\) 的梯度按式~\eqref{eq:potential-force} 诱导；执行器通道按式~\eqref{eq:actuator-dynamics} 以正定时间常数表达一阶滞后。这些性质在任意参数取值下都成立，因此构成不随训练改变的结构约束。

与之相对，逆质量的具体数值、耗散与耦合的方向依赖、广义力端口对相对速度与海流上下文的条件化方式，以及势能面的形状，都是数据驱动的可训练自由度。训练只在上述保结构的函数类内调整这些自由度，使长期递推轨迹拟合受控仿真数据。为使结构对照不被网络容量差异混淆，完整结构化模型、结构消融与黑箱基线采用容量匹配的网络宽度；海流任务统一以 REMUS~100 的质量矩阵与执行器先验初始化（第~\ref{subsec:candidates-evidence} 小节）。由此，模型族内的性能差异可在容量可比的前提下，归因于结构配置而非参数规模。

\subsection{控制块训练目标}
\label{subsec:training-objective}

训练目标是实验设置的组成部分，规定模型\emph{如何}在受控仿真数据上拟合，而不预设任何性能结论。训练数据由控制块序列组成：在单个控制块内，控制命令与外源上下文（海流、深度）保持分段常值，模型从块初值出发积分生成预测轨迹。预测过程为
\begin{equation}
  y_0=\mathcal{T}_{d\to m}(s_0),\qquad
  \hat y(t_i)=\Phi_\theta(t_i;t_0,y_0),\qquad
  \hat s_i=\mathcal{T}_{m\to d}(\hat y(t_i)),
  \label{eq:block-prediction}
\end{equation}
其中 \(\mathcal{T}_{d\to m}\) 与 \(\mathcal{T}_{m\to d}\) 为式~\eqref{eq:data-to-model-state} 与式~\eqref{eq:model-to-data-state} 给出的数据态--模型态转换，\(\Phi_\theta\) 表示由式~\eqref{eq:complete-augmented-vector-field} 的向量场经神经常微分方程积分得到的流映射。训练时，目标轨迹也转换到模型态 \(y_i=\mathcal{T}_{d\to m}(s_i)\)，使有海流样本的速度损失作用于相对水速度表示 \(\nu_r\)（依式~\eqref{eq:velocity-contract} 的速度契约）；位置、姿态与执行器状态在数据态与模型态中保持同一数值。一个典型加权训练目标可写为
\begin{equation}
  \mathcal{L}
  =
  \sum_i
  \lambda_x\norm{\hat x_i-x_i}^2
  +\lambda_R d_{\SO(3)}(\hat R_i,R_i)^2
  +\lambda_\nu\norm{\hat\nu_{r,i}-\nu_{r,i}}^2
  +\lambda_a\norm{\hat u_{a,i}-u_{a,i}}^2
  +\lambda_{\SO}\mathcal{L}_{\SO(3)} ,
  \label{eq:training-loss}
\end{equation}
其中各权重 \(\lambda_\bullet\) 记录于运行配置。姿态误差以旋转测地距离 \(d_{\SO(3)}(\hat R,R)\) 度量、而非旋转矩阵的欧氏差，使监督直接作用于 \(\SO(3)\) 流形上的几何偏离，并对脱离旋转群的非物理偏移保持敏感；\(\mathcal{L}_{\SO(3)}\) 为旋转矩阵的正交性与行列式正则项，约束预测姿态在数值积分下保持接近旋转群。速度项以相对水速度 \(\nu_r\) 为监督对象，确保水动力建模变量与第~\ref{sec:state-current-contract} 节的速度契约一致。评估与报告时，预测轨迹再通过 \(\mathcal{T}_{m\to d}\) 回到数据态，因此终端位置、姿态与总体速度等可观测指标仍在数据空间解释（第~\ref{subsec:metrics} 小节）。

\subsection{长期自由递推验证协议}
\label{subsec:rollout-protocol}

块级训练目标只检验局部受控初值问题的拟合质量，不能直接说明模型在自由递推中的误差累积。本文因而以\emph{长期自由递推}作为主要验证手段：将控制块依时间串接，首个控制块由给定初值启动，其后每个控制块以上一块的预测末状态作为模型态初值，并按数据协议更新位置原点、海流速度、控制命令与深度上下文。所得轨迹不再由后续真实状态校正，从而把单块拟合误差沿时间逐块传递、放大，直接暴露姿态漂移、速度发散、深度越界、常微分方程求解失败与非有限预测等只在长期才显现的失败模式。

长期递推使用受控仿真生成的 AUV 轨迹，区分无海流与有海流设置；有海流数据继续按式~\eqref{eq:velocity-contract} 维持总体速度与相对水速度的转换，使位姿几何按总体速度推进、而水动力配对作用于相对水速度。评估时域取 \(10\)、\(30\) 与 \(60~\mathrm{s}\) 三层，以 \(60~\mathrm{s}\) 为本章主时域；较短时域用于刻画误差随时间的增长形态与提前失败的位置，避免仅以单一终端时刻掩盖中途已发散的轨迹。一条轨迹若未完成目标时域，其终端误差\emph{不}计入该时域的条件误差统计，而由完成率与失败原因反映可比性。失败原因区分五类：真值轨迹本身失败、模型预测发散、常微分方程求解失败、非有限预测，以及深度或速度约束违反。这一区分使后续证据能把"模型结构导致的长期发散"与"数值或数据侧的局部失败"分离开来。

\subsection{评价指标与内部诊断量}
\label{subsec:metrics}

外部任务指标在数据空间定义，刻画模型作为长期轨迹预测器的可观测性能。主精度口径取 \(60~\mathrm{s}\) 终端位置误差的中位数与第 95 百分位（P95），二者\emph{并列}为主指标，分别刻画典型误差与尾部风险，不以单一标量排序——这一选择在第~\ref{sec:results-structure-evidence} 节将显示其必要性，因为部分模型的中位数与尾部排序并不一致。完成目标时域的轨迹另报告深度误差、旋转测地误差与总体线速度、角速度误差，并可给出对应的均值、最大值或随时域增长的误差曲线。完成率用于判断评估样本是否足够健康，只作为健全性指标，不直接参与精度排序：只有在完成目标时域的样本上，终端误差才具可比性。

内部诊断量在模型空间定义，用于说明各结构组件是否按设计工作，但不替代数据空间的外部任务误差。相对水速度误差用于核验水动力建模变量 \(\nu_r\) 是否与第~\ref{sec:state-current-contract} 节的速度契约一致；机械能量跨度与最大能量变化用于说明机械存储是否在递推中保持良定、有限的量级；\(\SO(3)\) 行列式与正交性误差用于说明姿态矩阵在长期递推中是否保持接近旋转群。聚合口径统一为：先在单个运行内对评估轨迹取统计量（如终端位置误差的中位数与 P95），再在随机种子之间对这些运行级统计量取均值；记号 \(N\) 表示异常排除后参与聚合的有效种子数（异常判据见第~\ref{subsec:candidates-evidence} 小节）。

\subsection{协议轴与初值扰动设计}
\label{subsec:protocol-axes}

初值扰动作为鲁棒性设计，分别施加于训练与评估两个阶段；二者是不同的干预，不应合并为同一种“协议敏感性”。为此本章把初值扰动沿三条\emph{正交}的轴组织，如表~\ref{tab:protocol-axes} 所示。训练协议轴（\texttt{train\_protocol}）刻画训练阶段是否、以及如何对块初值加噪；评估配置轴（\texttt{eval\_profile}）刻画评估阶段初值扰动的强度与偏置类型；评估协议轴（\texttt{eval\_protocol}）刻画评估端生成被扰动初值的采样方式。同一种轻量初值扰动既可作为训练线，也可作为评估端采样实现，其含义由所处的轴决定，不可互相替代或混为一谈。

\begin{table}[htbp]
  \centering
  \caption{初值扰动的三条正交协议轴}
  \label{tab:protocol-axes}
  \small
  \begin{tabularx}{\textwidth}{>{\raggedright\arraybackslash}p{0.20\textwidth}>{\raggedright\arraybackslash}p{0.24\textwidth}>{\raggedright\arraybackslash}X}
    \toprule
    协议轴 & 取值 & 含义 \\
    \midrule
    训练协议 \texttt{train\_protocol} & 干净；逐控制块独立噪声；轻量协议变体 & 训练阶段初值扰动的实现结构：干净初值；每个控制块独立采样噪声初值；或先生成轨迹级噪声观测、再令同一轨迹各控制块共享同一次噪声实现。 \\
    评估配置 \texttt{eval\_profile} & 干净；名义；退化；航向偏置 & 评估阶段初值扰动的强度与偏置类型，从干净到递增强度，以及叠加固定航向偏置的有偏配置；海流设置下另可附加海流估计偏置配置。 \\
    评估协议 \texttt{eval\_protocol} & 干净；逐样本独立；轨迹一致 & 评估端生成被扰动初值的采样方式，与训练协议轴的实现结构对应但作用于评估阶段。 \\
    \bottomrule
  \end{tabularx}
\end{table}

初值扰动的具体噪声模型遵循\emph{仅对初值加噪、基于 profile、在常微分方程语义上一致}的原则。噪声只作用于 \(t=0\) 的初始状态：从干净数据态出发转为干净模型态后，在模型态的相对水速度 \(\nu_r\)、姿态 \(R\)、执行器反馈 \(u_a\) 与海流 \(v_c\) 等通道上采样零均值扰动，得到被扰动的模型态初值再启动递推。这样做控制的是模型真正消费的变量误差预算，并在海流设置下依式~\eqref{eq:velocity-contract} 保持 \(R\)、\(\nu_r\) 与 \(v_c\) 的语义一致，避免“数据态看似扰动不大、而模型态实际输入已被过度污染”的失配。需要强调，本章的噪声接口定位为\emph{滤波后状态的鲁棒性正则}，而非完整导航误差后验：各通道误差预算主要由 profile 常数与训练集统计量确定，重点是控制模型消费变量的误差半径，而不是复现导航滤波器的联合协方差结构。

各通道的 profile 预算以 REMUS~100 的航位推算与惯性量级为参照确定，列于表~\ref{tab:noise-budget}。其中相对速度噪声按通道自然尺度设定为 \(\sigma_i=\max(\text{floor}_i,\ \alpha\,\mathrm{std}_i)\)，\(\mathrm{std}_i\) 取自训练集统计、\(\alpha\) 由 profile 决定、下限 \(\text{floor}_i\) 保证噪声不至不现实地趋零；姿态误差为 yaw 主导的各向异性小角度扰动，经 \(\SO(3)\) 指数映射作用到旋转矩阵。由此带来一个口径上的限定：名义与退化配置的部分语义是\emph{相对训练数据分布}的（经 \(\alpha\,\mathrm{std}_i\) 定义），而非完全固定的绝对物理规格；因此同一数据集内的跨模型比较是合理的，但 profile 的实际物理强度在跨数据集时可能漂移。为此，各运行在产物中额外记录各通道最终落地的 \(\sigma_i\)、姿态逐轴标准差与偏置幅值，使 profile 强度可追溯而不被误读为绝对不变量。

\begin{table}[htbp]
  \centering
  \caption{初值扰动各通道的 profile 噪声预算（REMUS~100 航位推算/惯性量级参照）}
  \label{tab:noise-budget}
  \small
  \begin{tabularx}{\textwidth}{>{\raggedright\arraybackslash}X cccc}
    \toprule
    通道（参数） & 名义训练 & 名义评估 & 退化评估 & 航向偏置评估 \\
    \midrule
    相对线速度 \(\nu_r\)（\(\alpha\)，下限 \(0.005~\mathrm{m/s}\)）& 0.03 & 0.05 & 0.10 & 同名义 \\
    相对角速度 \(\nu_r\)（\(\alpha\)，下限 \(0.0015~\mathrm{rad/s}\)）& 0.03 & 0.05 & 0.10 & 同名义 \\
    姿态（yaw 主导小角，基础量级 / rad）& 0.0035 & 0.0050 & 0.0120 & 名义 \(+\) 固定 0.015 偏置 \\
    海流 \(v_{cx},v_{cy}\) / \(v_{cz}\)（m/s）& 0.008 / 0.004 & 0.012 / 0.006 & 0.030 / 0.015 & 同名义 \\
    执行器 \(\delta_r,\delta_s\)（rad）/ 转速（rpm）& 0.002 / 3 & 0.003 / 5 & 0.008 / 15 & 同名义 \\
    \bottomrule
  \end{tabularx}
  \par\vspace{2pt}
  \footnotesize 注：相对速度噪声为 \(\sigma_i=\max(\text{floor}_i,\alpha\,\mathrm{std}_i)\)，表中给出 \(\alpha\) 与下限；姿态为 yaw 主导各向异性小角度经 \(\SO(3)\) 指数映射。航向偏置评估在名义评估的随机扰动上叠加固定幅值 \(0.015~\mathrm{rad}\) 的 yaw 偏置，偏置符号按样本固定、在一条递推内不随时间变化。名义与退化的部分语义相对训练集分布定义，各运行记录实际落地预算。
\end{table}

训练端采用渐进噪声调度以降低噪声初值直接打崩优化的风险：默认前 \(20\) 个 epoch 完全使用干净初值，随后 \(80\) 个 epoch 将噪声强度从零线性爬坡至目标值，达到稳态后约半数训练样本使用噪声初值、其余保持干净（以样本级掩码实现）。评估端则对所有模型固定噪声随机种子，以保证不同模型之间的被扰动初值可直接比较。

\subsection{候选模型、结构消融与证据纳入规则}
\label{subsec:candidates-evidence}

\paragraph{候选比较空间。}候选比较空间由本文的可训练模型族确定，用于在同一状态契约、训练损失与长期递推评估口径下形成结构对照，而非预设性能排序。比较链条覆盖完整结构化模型、端口哈密顿风格基线、\(\SE(3)\) 加速度与动量黑箱、全状态黑箱，以及完整模型内部的结构消融，对应显式位姿几何、动量与质量映射、机械存储与广义力组织、耗散与零功率耦合四类检验维度。其中 \(\SE(3)\) 几何与动量基线、端口哈密顿风格基线参照 \(\SE(3)\) 哈密顿及李群端口哈密顿神经常微分方程的建模思路~\citep{duong2021se3hamnode,duong2024liegroupphnode}，但均在统一的 AUV 状态契约、数据协议与评估口径下重新实现为受控候选，而非直接复刻文献中的原始模型与训练设置。表~\ref{tab:model-comparison-roles} 汇总各模型类别在比较空间中的角色。

\begin{table}[htbp]
  \centering
  \caption{候选模型比较与结构消融的检验目标}
  \label{tab:model-comparison-roles}
  \begin{tabularx}{\textwidth}{>{\raggedright\arraybackslash}p{0.26\textwidth}>{\raggedright\arraybackslash}p{0.31\textwidth}>{\raggedright\arraybackslash}X}
    \toprule
    模型类别 & 代表模型 & 检验目标 \\
    \midrule
    完整结构化模型 & AUVHamNODE & 同时包含速度契约、\(\SE(3)\) 运动学、机械存储、\(D_\theta\)/\(J_\theta\)/\(\tau_\theta\) 分解和执行器滞后，作为本文结构化假设空间的完整候选项。 \\
    端口哈密顿风格基线 & 合并非保守广义力基线、位形广义力基线 & 比较机械存储核心保留时，不同非保守广义力组织方式对长期递推行为的影响。 \\
    \(\SE(3)\) 加速度黑箱 & 显式 \(\SE(3)\) 运动学与黑箱加速度动力学 & 检验保留位姿几何但不保留完整机械存储分解时的候选对照。 \\
    \(\SE(3)\) 动量黑箱 & 显式 \(\SE(3)\) 运动学、常质量映射与黑箱动量动力学 & 检验动量坐标和质量映射保留时，黑箱动量动力学的候选对照。 \\
    全状态黑箱 & 非结构化状态导数模型 & 提供缺少显式位姿几何和机械能量结构时的长期递推风险对照。 \\
    结构消融 & 无质量先验、对角耗散、无零功率升力耦合、广义力仅执行器条件化 & 提供质量先验、耦合耗散、斜对称零功率耦合和开放广义力条件化范围的逐项结构对照。 \\
    \bottomrule
  \end{tabularx}
\end{table}

\paragraph{结构消融。}结构消融在完整模型上逐项移除或弱化某一结构组件，用于评估第~\ref{sec:energy-power-properties} 节中各结构配置对长期递推行为的影响。每个消融对象对应一个明确的候选结构假设：

\begin{itemize}
  \item 质量先验消融移除或弱化基于物理质量矩阵的先验初始化，同时保留正定质量参数化，用于评估质量先验对动量--速度配对学习和长期递推行为的影响。
  \item 对角阻尼消融将半正定耦合耗散矩阵 \(D_\theta(\xi)\) 限制为非负对角形式，用于评估跨通道耗散耦合相对独立通道耗散的作用。
  \item 零功率耦合消融移除斜对称分支 \(J_\theta(\xi)\nu_r\)，用于评估升力或横向水动力类零功率耦合在候选模型族内的作用；该消融不等同于证明真实水动力升力项被唯一辨识。
  \item 广义力条件化消融把广义力分支 \(\tau_\theta\) 收缩为主要由执行器状态或命令驱动，弱化相对速度和海流上下文条件化，用于评估开放广义力通道的条件化范围对长期递推的影响。
\end{itemize}

\paragraph{统一实验配置。}为保证长期递推结果可复现，所有候选模型在统一的数据、训练与数值实现配置下比较\footnote{连续时间积分基于 torchdiffeq~\citep{chen2018neuralode} 的固定步长四阶 Runge--Kutta 实现。}，完整的命令级参数、噪声预算与运行元数据保留在各运行的配置、评估摘要与溯源记录中。表~\ref{tab:experimental-config} 汇总与论证相关的主要设置。

\begin{table}[htbp]
  \centering
  \caption{长期递推实验的主要配置}
  \label{tab:experimental-config}
  \begin{tabularx}{\textwidth}{>{\raggedright\arraybackslash}p{0.18\textwidth}>{\raggedright\arraybackslash}X}
    \toprule
    配置维度 & 设置 \\
    \midrule
    数据规模 & 海流数据集含一千条轨迹、每条一百五十个控制块；控制块时长 0.2~s，状态采样间隔 0.05~s，训练与验证按 \(0.8\) 划分。 \\
    速度契约 & 数据态保存总体速度，训练前按式~\eqref{eq:velocity-contract} 转换为模型态相对水速度。 \\
    数值积分 & 固定步长四阶 Runge--Kutta，单精度浮点状态；训练、块级评估与长期递推共用同一求解器与损失口径。 \\
    网络容量 & 完整结构模型、结构消融与黑箱基线采用容量匹配的网络宽度；海流任务以 REMUS 质量矩阵与执行器先验初始化。 \\
    优化 & AdamW，学习率经线性预热后余弦衰减，至多三百个 epoch；批量、学习率端点与正则权重记录于运行配置。 \\
    递推评估 & 重采样初值，覆盖 PRBS、CHIRP 与 OU 三类控制场景，每场景 30 条评估轨迹；评估时域取 10、30 与 60~s，主时域 60~s。 \\
    \bottomrule
  \end{tabularx}
\end{table}

\paragraph{证据纳入与异常判据。}进入定量比较的运行须满足统一的数据协议、状态转换、评估口径，以及完整的随机种子、程序版本与运行环境记录。本章对运行级异常区分两类相互独立、触发后均移出定量聚合并单独披露的判据。其一为\emph{训练异常}：训练阶段出现灾难性梯度失败（如训练日志出现“无成功训练批次”特征且计数大于零）。其二为\emph{递推发散}：训练收敛，但长期递推关键终端指标缺失、非有限或超过预设发散阈值。两类异常以训练异常优先计入、同一运行不重复计数；移出定量聚合后按类别与随机种子单独披露，以免选择性举证。若某模型在同一协议单元下全部种子均递推发散，则不报告有限中位数，而记为\emph{长期稳定性失败}。当不同代码环境或数据镜像的证据相互冲突时，以数据协议一致性、记录完整性与失败可复现性判定统一证据状态，而非按来源一概取舍；证据状态的完整判定规则见溯源与数据目录文档。需说明，本节只规定候选比较空间、消融设计与纳入规则；实际进入定量比较的证据子集——即哪些 \((\text{模型}\times\text{协议}\times\text{种子})\) 组合在统一代码环境与数据镜像下被重建并纳入——在第~\ref{sec:results-structure-evidence} 节的当前证据口径下报告，本节不预先锁定。
